\section{AegisKV Design}
\label{sec:design}

AegisKV combines workflow delay tolerance with each request's KV
recovery demand. The Tracker estimates slack. The Router selects execution
and recovery DPs at arrival. The KV Master chooses backup tiers,
per-request layer depths, and hot blocks at backup-policy events.
The checkpoint interval is configured, while Adapters update target-local
Peer budgets at window boundaries. The following design
separates these decisions from the evidence that motivates them.

\subsection{Design Insights}
\label{sec:design-insights}

Three premises motivate the design (Figure~\ref{fig:design-insights}).
The first draws on published agent-prediction results; the other two
characterize selected sparse-attention releases and recovery demand.

\begin{figure*}[t]
  \centering
  \captionsetup[subfigure]{font=normalsize,labelfont=normalfont,
    textfont=normalfont,justification=centering,singlelinecheck=false,skip=3pt}
  \begin{subfigure}[t]{0.32\textwidth}
    \vspace{0pt}
    \includegraphics[width=\linewidth]{figures/agent-prediction.pdf}
    \caption{Agent trajectory prediction}
    \label{fig:insight-agent-time}
  \end{subfigure}\hfill
  \begin{subfigure}[t]{0.32\textwidth}
    \vspace{0pt}
    \includegraphics[width=\linewidth]{figures/sparse-models.pdf}
    \caption{Selected sparse releases}
    \label{fig:insight-sparse-adoption}
  \end{subfigure}\hfill
  \begin{subfigure}[t]{0.32\textwidth}
    \vspace{0pt}
    \includegraphics[width=\linewidth]{figures/retry-kv-prediction.pdf}
    \caption{Retry KV prediction}
    \label{fig:insight-retry-kv}
  \end{subfigure}
  \caption{Design evidence. (a) Adapted from Heddle~\cite{zhang2026heddle},
    Figure 13: Pearson correlation of predicted and actual remaining
    trajectory length after one and two agent steps, with Qwen3-14B.
    Values are extracted from the published vector figure; these are
    literature results, not AegisKV runtime-error measurements.
    (b) Selected releases and
    native context: 1 MiniCPM4, 2 MiniCPM4.1, 3 DeepSeek-V3.2-Exp,
    4 DeepSeek-V3.2, 5 GLM-5, 6 MiniCPM-SALA, 7 DeepSeek-V4 preview,
    8 MiniMax-M3, 9 GLM-5.2, 10 LongCat-Flash-Lite-Sparse.
    Dates follow announcements/model cards. Gray steps trace the sample
    maximum; markers denote top-$k$ units in sparse branches.
    (c) Three native PG19 captures/model: 32K history; queries
    32,257--32,768 predict distinct historical KV for 32,769--32,832.
    Objects are layer/64-token blocks containing all KV heads; both
    policies include mandatory blocks in the 40\% budget.
    Seer uses an experimental Qwen3-8B Prefill gate~\cite{song2026seerqwen3}
    (threshold $10^{-3}$; final-block dense override disabled).
    Bars/whiskers are document means/ranges, not fault-recovery runs.}
  \label{fig:design-insights}
  \Description{Three side-by-side panels show Heddle's published agent
    trajectory prediction correlations, a primary-source survey of sparse model releases,
    and held-out retry KV coverage. Past-frequency ranking improves
    coverage substantially for NOSA and only modestly for MiniCPM4
    and SeerAttention-Qwen3.}
\end{figure*}

\paragraph{Insight 1: execution context predicts remaining agent work.}
Existing predictors exploit task semantics and execution progress.
Heddle~\cite{zhang2026heddle} reports trajectory-length correlations of
approximately 0.72, 0.88, and 0.94 for coding, search, and math agents
after two steps, improving over one-step predictions
(Figure~\ref{fig:design-insights}(\subref{fig:insight-agent-time})).
AgentDecarbonizer~\cite{yan2026agentdecarbonizer} additionally estimates
task durations and updates remaining time from observed progress.
AegisKV builds on these capabilities: the Tracker converts remaining-time
estimates into workflow slack for KV protection. Trajectory correlation
supports predictability; runtime calibration remains deployment-specific.

\paragraph{Insight 2: sparse attention changes the recovery objects.}
Long-context models use sparse mechanisms including InfLLMv2,
DSA, or compressed/hybrid attention~\cite{minicpm4model,
openbmb2025minicpm41,deepseek2025v32exp,deepseek2025v32model,
zai2026glm5model,openbmb2026sala,deepseek2026v4model,minimax2026m3,
zai2026glm52,longcat2026sparse}.
Figure~\ref{fig:design-insights}(\subref{fig:insight-sparse-adoption})
establishes release examples, not deployment share.
The KV Master therefore consumes actual layer/block selections and
model-specific recovery state instead of valuing all KV bytes equally.

\paragraph{Insight 3: past selections predict part of retry demand.}
With 40\% of each model's historical KV bytes, historical frequency
covers 77.8\%, 53.6\%, and 43.0\% of distinct needed KV for NOSA-3B,
MiniCPM4-8B, and SeerAttention-Qwen3-8B, versus 59.7\%, 50.9\%, and
42.0\% for recency
(Figure~\ref{fig:design-insights}(\subref{fig:insight-retry-kv})).
Seer's demand spans 92.5\% of historical blocks after the head union,
limiting even oracle placement to 43.3\% coverage. Thus gains depend on
the model and placement granularity. These held-out book-text suffixes
motivate advance placement with exact demand fetches for misses;
they do not establish performance on agentic failure retries.
\subsection{System Overview}
\label{sec:system-overview}

Figure~\ref{fig:system-overview} separates applications, control, and
inference engines. Applications report request/workflow IDs and DAG
updates. The Router uses load and KV locality for normal dispatch.
It also binds the recovery target used for Peer protection. The Controller maintains
slack, backup policy, and surviving directory and route metadata.
KV Master sets each request's depth and hot blocks at checkpoints
or policy changes. Each Adapter manages its local Peer budget
and invokes a Connector to copy KV directly
between engines and storage. Badges index the mechanisms below, not
execution order.

\begin{figure}[t]
  \centering
  % User-requested single-column placement; preserve the source aspect ratio.
  \includegraphics[width=\columnwidth]{figures/system-overview.pdf}
  \caption{AegisKV overview. Badges identify Tracker (1), recovery-aware
    Router (2), KV Master (3), and Adapter (4). Router dispatches normal
    requests and binds retry/Peer targets at arrival. The dispatch and
    retry arrows illustrate a DP0 request taken over by DP1.
    KV Master selects depth and hot blocks; Adapters control Peer budgets.
    Peer and DRAM paths commit independently.
    Equivalent DPs hold local KV and smaller Peer regions; memory
    proportions are schematic.}
  \label{fig:system-overview}
  \Description{Three horizontal tiers show applications, control, and inference.
    A workflow DAG feeds a Tracker with per-request slack bars. One
    recovery-aware Router selects normal execution DPs, binds retry and
    Peer targets, and activates retries after confirmed failure. The KV Master
    selects protection tiers, request-specific depths, and hot-block candidates.
    Equivalent DP0 and DP1 engines each contain an Adapter controlling Peer budgets,
    a layer/block grid in Local KV, a smaller Peer KV region holding the other
    DP's replicas, and a Connector. Adapters invoke Connectors; Peer transfers
    and independent DRAM copy/restore paths bypass the control plane.}
\end{figure}

The Router binds retry and Peer targets alongside normal dispatch
(\S\ref{sec:recovery-router}). KV Master uses a fixed checkpoint
interval and reviews each request's depth at policy events
(\S\ref{sec:kv-master}). Adapters update budgets at window boundaries
(\S\ref{sec:kv-runtime}). The Controller
retains directory and commit metadata outside the source's failure domain.
\subsection{Workflow Tracker: GBDT and CPM}
\label{sec:workflow-protection}

Tracker (1) predicts remaining execution time with GBDT, then computes
slack by CPM~\cite{kelley1959cpm} (Algorithm~\ref{alg:slack-estimation}).
GBDT is a replaceable predictor instantiation; the literature evidence
in Figure~\ref{fig:insight-agent-time} concerns Heddle's predictor.
At time $t$, the residual DAG $G_t=(V_t,E_t)$ contains unfinished agent
execution segments. A segment groups sequential requests between
dependency barriers; parent agents split at fork/join boundaries so
child execution is not counted again as parent waiting. Completed
segments are removed, and each live request $r$ maps to its segment $a(r)$.

\paragraph{Predict remaining time.}
The GBDT $f_\theta$ uses role/model, elapsed time, known prompt lengths,
and completed-request counts, output lengths, and latencies.
Historical running snapshots have target $t_a^{\mathrm{finish}}-t$;
entry snapshots use segment duration, masking metadata unavailable
before launch. Thus unstarted segments predict execution after their
dependencies become ready. Future requests, outputs, and completion
times are never input features. GBDT predicts the remaining duration
directly; elapsed time is not subtracted again.

\begin{figure}[t]
  \raggedright
  \refstepcounter{designalgorithm}\label{alg:slack-estimation}
  \noindent\textbf{Algorithm \thedesignalgorithm: GBDT--CPM slack estimation}
  \par\smallskip
  \noindent Input: $G_t$, visible features $x(t)$, GBDT $f_\theta$, live
  requests $\mathcal R_t$, and segment map $a(r)$.
  Output: request slack $\sigma$.
  \par\smallskip
  \begin{tabularx}{\columnwidth}{@{}r@{\quad}X@{}}
    \toprule
     1 & $\sigma_r\gets0$ for all $r\in\mathcal R_t$ \\
     2 & \textbf{if} DAG structure incomplete: \textbf{return} $\sigma$ \\
     3 & $z_a\gets f_\theta(x_a(t))$ for all $a\in V_t$ \\
     4 & \textbf{if} any prediction unavailable or invalid: \textbf{return} $\sigma$ \\
     5 & $\widehat d_a\gets\max(\varepsilon,z_a)\ \forall a$; $\pi\gets\mathrm{Topo}(G_t)$ \\
     6 & \textbf{for} $a$ in $\pi$: \\
     7 & \quad $\mathrm{EF}_a\gets\widehat d_a+\max_{b\in P(a)}\mathrm{EF}_b$ \\
     8 & $T\gets\max(\{0\}\cup\{\mathrm{EF}_a:a\in V_t\})$ \\
     9 & \textbf{for} $a$ in reverse $\pi$: \\
    10 & \quad $\mathrm{LF}_a\gets\min_{b\in S(a)}(\mathrm{LF}_b-\widehat d_b)$ \\
    11 & \textbf{for} $r\in\mathcal R_t$: \\
    12 & \quad $\sigma_r\gets\max(0,\mathrm{LF}_{a(r)}-\mathrm{EF}_{a(r)})$ \\
    13 & \textbf{return} $\sigma$ \\
    \bottomrule
  \end{tabularx}
  \par\smallskip
  \noindent $P(a)$/$S(a)$: predecessors/successors. Empty predecessor
  max is $0$; empty successor min is $T$; $\varepsilon>0$ is a
  minimum remaining-time floor.
  \Description{GBDT predicts remaining agent-segment durations. Forward
    CPM computes earliest finishes and the workflow horizon; backward
    CPM computes latest finishes. Requests inherit their segment's slack.
    Incomplete dependency structure or unavailable predictions return zero slack.}
\end{figure}

\paragraph{Propagate and refresh slack.}
The forward pass computes earliest finishes relative to now; the
backward pass computes latest finishes without extending horizon $T$.
Their difference is the segment slack inherited by its live requests.
This shared margin is not an independent allowance for every request.
Request arrivals/completions, DAG updates, and placement windows trigger
refreshes; CPM costs $O(|V_t|+|E_t|)$. Unknown dependency structure or missing
predictions force zero slack for that workflow. These are planning
estimates, not delay guarantees; system evaluation must account for
the selected predictor's runtime errors and overhead.

The orchestrator must declare the remaining fork/join structure complete
and map every live request to a segment. An unexpanded branch, unresolved
dependency, or missing mapping makes it incomplete; discovering new work
invalidates the prior declaration and triggers refresh. Queueing features
may inform durations, but CPM does not model future resource contention.
Train/test splits must separate workflows and time periods, not snapshots
of one trajectory. Required diagnostics include seconds-scale error,
false-positive slack, critical-path ranking, zero-slack fallback frequency,
and decision regret against zero-slack, historical, and oracle baselines.
Published correlation alone validates none of these runtime properties.
\subsection{Recovery-aware Router}
\label{sec:parallel-recovery}
\label{sec:recovery-router}

Router (2) handles normal dispatch and recovery routing.
Its extension to normal routing is an arrival-time recovery assignment:
the DP that will receive a retry is also the destination of any
Peer-HBM protection selected later by the KV Master.

\paragraph{Dispatch normally and bind a recovery target.}
On arrival, the Router assigns request $r$ to a source DP $g_r$ using
foreground load and compatible KV locality. It also records a compatible
recovery DP $p_r$ outside the source's declared failure domain.
For $n\ge2$ equivalent, failure-separated DPs, a simple example is
$p_r=(g_r+1)\bmod n$; per-request assignments can spread recovery load
across several targets. This example is a binding policy, not a claim of
optimal failover load balance. A TP group is one DP and must restore
compatible shards together.

The Router publishes $(r,g_r,p_r)$ and its route version to surviving
control state and the KV Master. A request running on DP0 with target
DP1 therefore places any admitted Peer copies on DP1, where its retry
can use them locally. Binding reserves neither Peer capacity nor recovery
workspace. The Master decides the protection tier, layers, and blocks;
the target Adapter controls admission. The Router does not choose a
request's hot set or initiate an unconditional backup.

\paragraph{Activate the binding after a confirmed failure.}
Health monitoring combines membership, device errors, and GPU progress
reported by Adapters; a process heartbeat alone does not establish GPU
progress. Once the failure is confirmed and the old source is fenced,
the Router retrieves affected requests, route versions, acknowledged
token prefixes, and committed replica locations from surviving control
state, then submits retries to the recorded targets. Peer replicas are
local to the assigned target; DRAM loading or recomputation supplies the
remaining state through the engine and Connector.

The Router rechecks target health, model/layout compatibility, and
execution-space feasibility. If a target fails or cannot admit the
working set, the request waits or receives a new valid assignment.
Remapping updates the route version, informs the Master, and resolves
replicas again: a copy on the old target cannot be credited as local
HBM on the new one. A route change alone does not migrate or commit KV.

\paragraph{Consume a valid recovery state.}
The Router resolves the latest dependency-complete checkpoint manifest
from the directory; publication rules are described in
\S\ref{sec:kv-master}. Partial copies do not advance that frontier.
An incomplete new checkpoint leaves the older valid frontier usable;
without one, the engine reconstructs from recoverable inputs.

A surviving authority assigns a monotonically increasing source epoch;
copies, directory commits, and output acknowledgements reject stale epochs.
Before acknowledging a generated token, the runtime persists its request
ID, position, token ID, and epoch in a surviving ordered log. Prompt and
accepted tool inputs are also recoverable. KV may lag this log; recovery
reconstructs the acknowledged suffix and suppresses duplicate delivery.
External tool side effects require the orchestrator's own idempotency
contract. The Router uses these runtime guarantees to identify a valid
retry; it does not itself execute KV reconstruction or tool replay.

\paragraph{Submit retries without a completion barrier.}
Requests become eligible independently after state and target checks.
Slack orders protection decisions, not a rule that one retry must finish
before another starts. The engine admits feasible work and advances it
at supported batching and kernel boundaries while continuing foreground
service. The Adapter and Connector provide the resource and progress
interfaces in \S\ref{sec:kv-runtime}. Local backup does not itself provide execution
space or eliminate queueing. The routing and recovery interfaces above
are design contracts whose concurrent implementation requires validation.

\subsection{KV Master: Request-Specific Backup Decisions}
\label{sec:kv-placement}
\label{sec:kv-master}

KV Master (3) selects recovery tiers, each request's Peer depth $L_r$,
and layer-local hot sets. The Router supplies the recovery target;
the Adapter publishes its window-level capacity $B_p$. Requests may use
different depths within the same capacity window.

\paragraph{Decision events.}
The Master establishes policy at arrival, offers sealed state after prefill,
and reviews a request every additional $Z$ generated tokens. $Z$ is a
block-aligned system configuration, not an Adapter output. Significant
slack, budget, context, or copy-state changes can trigger an earlier review.
A decision does not imply a completed checkpoint: the replay tail extends
beyond the latest committed frontier and can exceed $Z$ when copying lags.
Copy queues have configured byte and operation limits. Coalesce superseded
unstarted plans, retain dependencies of in-flight or committed versions,
and skip a new backup attempt if it cannot be reserved. Report the resulting
commit lag; the trigger interval supplies no durability bound.
Sealing only updates eligibility. Hot-block statistics use the preceding
512 decode steps, independently of checkpoint and capacity windows.

\paragraph{Publish an incremental checkpoint.}
Publish a checkpoint manifest only after its full dependency closure is
committed: compatible model, weights and layout, all required layers and TP
shards, sealed token intervals, selector state, and the matching token-log
prefix. A frontier is the latest valid common manifest, not the maximum
of independently updated object locations. Peer and DRAM may commit in
either order; each manifest entry must resolve to a matching surviving
replica. An incomplete new manifest leaves the older one authoritative;
without one, replay from recoverable inputs. Recovery validates versions
and closure again after fencing. These checks define the minimum contract;
their concurrent implementation remains to be validated.

\paragraph{Choose protection tiers.}
Let $T_{\mathrm{rc}}$ denote prefix recomputation, $T_{\mathrm{rp}}$ known-tail
replay, $W_0$ exposed DRAM-loading wait, and $R_{\mathrm{floor}}$ detection,
rerouting, and readmission time. Joint context/tail/load profiles give
\begin{equation}
\begin{aligned}
 R_{\mathrm{base}}&=R_{\mathrm{floor}}+T_{\mathrm{rp}},\\
 R_R&=R_{\mathrm{base}}+T_{\mathrm{rc}},\qquad
 R_D=R_{\mathrm{base}}+W_0.
\end{aligned}
\label{eq:recovery-base}
\end{equation}
$R_R$ and $R_D$ use the weighted conditional means of the same
future-state samples defined below; the tier test is a planning heuristic,
not a deadline guarantee. Tier R needs no backup if $\sigma_r\ge R_R$.
Otherwise choose D when $R_D<R_R$ and $\sigma_r\ge R_D$; remaining
requests compete for Peer as tier P. There is no assumption that loading
is faster than recomputation. The non-Peer fallback chooses the cheaper
feasible route, preferring R on equal costs; a denied Peer increment does
not force a slower DRAM restore. DRAM protection requires its own admission
and commit (\S\ref{sec:kv-runtime}). Missing profiles request conservative
DRAM protection when feasible, without claiming it is faster; failed
admission uses an older valid checkpoint or R. Slack orders protection,
not retry execution. At recovery, refresh costs and availability; overlapping
cost intervals retain the feasible current route unless dominance is clear.

\paragraph{K counts resident blocks, not selector outputs.}
A protection object is an immutable sealed \emph{(layer, block)} including
all KV heads; the trace study uses 64-token blocks. Let $H_{r,l}$ be
request $r$'s published, committed Peer set at layer $l$. We define
\begin{equation}
 K_{r,l}=|H_{r,l}|.
 \label{eq:hot-block-count}
\end{equation}
Thus $K$ counts Peer-resident blocks of one request and layer. It is not
the model's per-query top-$k$, a hit-rate guarantee, or a count of copies
still in flight. Physical allocation also charges retiring buffers,
reservations, and selector state; shared copies are counted once per pool.

A separately configured quota determines the desired hot-set size.
The configuration declares either a fixed block limit or
$\lceil h_l n_r\rceil$ blocks for $n_r$ sealed historical blocks.
Mandatory sealed sink/window blocks are included in that limit: if they
alone exceed it, request a larger explicitly charged quota or decline
that layer's Peer admission, never silently drop native requirements.
The rule is
fixed in each comparison; this paper does not claim an optimal quota.
Remaining slots rank native selection counts, then recency and stable
block ID. Optional blocks leave after two consecutive statistics windows
outside the target set. Required auxiliary state is budgeted separately.
Before 512 steps are available, use the observed prefix; with no observations,
rank mandatory blocks first and then recency and stable ID. This cold-start
rule has no demonstrated prediction accuracy.
For a replay of $N$ tokens, demand is the deduplicated union of native
selections across queries and KV groups, not $N$ copies of a single top-$k$.

\paragraph{Preview replacement as well as admission.}
A candidate depth $d$ offers target hot sets in layers $1{:}d$;
this prefix family restricts the search, rather than proving that arbitrary
layer placements cannot do better.
Its plan retains intersections, demotes owned blocks leaving the set,
and admits new blocks. It cannot evict another request's protected copy.
A shared object loses only the departing owner's reference.
Dirty demotions, new transfers, and temporary space all enter the preview;
slots become reusable only after actual release
(\S\ref{sec:kv-runtime}). Shrinking $L_r$ can therefore retire the owner's
deeper copies. An unchanged depth can still replace block identities.
Uniform, fully admitted layers contain $L_r K$ layer-block objects;
partial admission is reported through the actual sets $H_{r,l}$.

\paragraph{A reproducible depth objective.}
All candidates start from the same target snapshot $S$ and use the same
planning horizon $H$, future-state samples $j$, and normalized weights $w_j$,
conditional on a source failure within $H$. Current slack is held fixed
within this local comparison and refreshed at the next policy event.
A sample specifies committed history $M_j$, replay tail $N_j$, available
copy time, and foreground/retry load. Samples come from a declared,
versioned profile, including commit lag; they do not assume $N_j=Z$ or
a uniform checkpoint phase. Co-located requests share the same source-event
samples and contended target resources. For a configured Poisson model,
$\pi_r=1-\exp(-\lambda_r H)$. Other risk models must provide the corresponding
probability over the same $H$; a failure rate is not itself a probability.

Let $G_d$ be the admitted-and-retired plan from the preview.
$A_r(d)$ is the predicted aggregate change in foreground request-delay
seconds, relative to keeping the current placement at $S$. The cohort
contains active requests and arrivals within $H$, followed to completion
under matched arrivals and a held policy, including completions after $H$.
It includes replication, writeback, transient occupancy, and reduced or
increased admission pressure, and may be negative. Let $R_r(d,j)$ be
failure-to-ready latency for sample $j$, crediting only replicas predicted
to commit before its failure. The Master minimizes
\begin{equation}
 J_r(d)=A_r(d)+\pi_r\sum_j w_j
          \max\{0,R_r(d,j)-\sigma_r\}.
 \label{eq:depth-score}
\end{equation}
Both terms have units of delay seconds. This is a local protection
heuristic in slack order; it neither sums to workflow makespan nor
optimizes all later requests' allocations.
For example, two equally critical branches on one failed source each
need 10\,s to recover. Protecting only one reduces their summed local
penalties from 20 to 10\,s but leaves workflow delay at 10\,s. Source-event
DAG replay and a feasible small-DAG oracle are required to assess this gap.

\paragraph{Profile interface and validity.}
A profile query takes the model/selector and weight revisions, KV layout,
block size, hardware/link topology, $S$, $G_d$, the shared state samples,
and the configured copy scheduler. It returns $A_r(d)$ and $R_r(d,j)$
with uncertainty intervals, a generation ID, support ranges, and expiry.
$A$ is obtained from matched normal-service traces with identical arrivals;
$R$ uses matched recovery traces or a declared, calibrated event model.
Transfer and compute overlap are evaluated on one timeline; adding their
full durations would double-count hidden transfer. Rows state whether
they are measured or modeled. Missing measurements are not filled with
the Figure~\ref{fig:design-pipeline} example.

The finite depth grid includes zero and the current depth, and an explicit
\emph{keep} plan preserves current block identities. The lowest score wins;
ties prefer smaller depth, then stable plan ID. Propagating the profile's
intervals through (\ref{eq:depth-score}) gives lower and upper score bounds.
A change is accepted only if its upper bound plus a configured margin
is below the keep plan's lower bound. Unsupported candidates are excluded.
An invalid generation or missing, expired, or drift-flagged keep profile
freezes discretionary expansion and replacement; new requests attempt
DRAM admission, with the stated fallback if it fails.
Existing protection is retained only while feasible:
Adapter contraction and correctness-driven retirement still take precedence.
A generation is invalidated after a configured number of consecutive
relative-error violations on comparable observed operations; unmatched
observations do not count as validation. Thresholds, grid, horizon,
risk model, quota rule, sample weights, and calibration version are
recorded configuration, not inferred experimental constants.
At zero failure probability, only $A$ affects the score; costly expansion
has no recovery justification, while a zero-cost keep plan need not move.
Uncertain route crossings do not justify an unprofiled migration.

\begin{figure}[t]
  \raggedright
  \refstepcounter{designalgorithm}\label{alg:placement}
  \noindent\textbf{Algorithm \thedesignalgorithm: Per-request backup policy}
  \par\smallskip
  \noindent Input: requests, slack, fixed $Z$, profiles, targets, budgets.
  Output: admitted tiers, Peer plans, and committed $H_{r,l}$.
  \par\smallskip
  \begin{tabularx}{\columnwidth}{@{}r@{\quad}X@{}}
    \toprule
    1 & Snapshot copies, leases, budgets, and reservations as $S$ \\
    2 & \textbf{for} affected $r$ in $(\sigma_r,\mathrm{ID}_r)$ order: \\
    3 & \quad Choose R/D/P and the feasible non-Peer fallback \\
    4 & \quad For D/P, request DRAM; report coverage/failure \\
    5 & \quad Rank layer-local candidates under the quota rule \\
    6 & \quad Form keep and depth plans; R/D target depth is zero \\
    7 & \quad Preview each plan independently from the same $S$ \\
    8 & \quad Query valid profiles; score by (\ref{eq:depth-score}) \\
    9 & \quad Choose a stable improvement, or feasible fallback \\
    10 & \quad Queue its retains, demotions, and new admissions \\
    11 & \quad Update $S$ with only this chosen plan \\
    12 & Atomically recheck reservations before physical execution \\
    13 & Publish only committed residency and actual released bytes \\
    \bottomrule
  \end{tabularx}
  \Description{The Master chooses per-request layers and hot sets using
    a shared-sample cost objective. Only chosen plans change the virtual
    allocation snapshot; execution rechecks space and safe retirement.}
\end{figure}

The preview is not a reservation. Actual grants may be smaller. The
directory reports committed $K_{r,l}$ and frontiers; a mismatch triggers
another review. Sorting $n$ requests takes $O(n\log n)$ time. Candidate
work scales with the grid, sample count, and affected blocks. Profile
calibration and accuracy remain evaluation work.

\paragraph{Keep selection and placement separate.}
Recovery restores model-specific selector state before invoking the
unchanged native selector. The selector may name any causally eligible
historical or recomputed KV, including DRAM-only objects. The Connector
resolves logical IDs to committed replicas, fetches actual misses, and
returns readiness events; attention waits for every required input.
Forecast mistakes add traffic or waiting, not a new attention approximation.
Tier R and missing committed prefixes use dependency-correct recomputation
of logged tokens. A batched replay is eligible only if its native selector
and state evolution are equivalent to the original token-wise path under
the declared numerical contract; otherwise replay token by token. A native
prefill gate or an unresolved numerical diagnostic does not establish this
equivalence. Compatible weights, layout, and operators are required;
future stochastic samples need not be bitwise identical across floating-point
execution schedules.

\begin{figure}[t]
  \centering
  \includegraphics[width=\columnwidth]{figures/recovery-pipeline.pdf}
  \caption{Modeled recovery pipelines: MiniCPM4-8B, 32 layers, 32K
    history, 64-token known tail; $C=2.87$\,ms/layer (A100-derived average)
    and $D=5.37$\,ms at an assumed 6.25\,GB/s. In (d), 147\,MiB of Peer
    residency removes preload wait. Totals end at tail-KV reconstruction.
    The homogeneous, preemptible-copy model excludes detection,
    queuing, and interference.}
  \label{fig:design-pipeline}
  \Description{Four timelines compare token-block transfer, full-layer
    pipelining, hot preloads with demand fetch, and Peer-resident hot sets.}
\end{figure}

Figure~\ref{fig:design-pipeline} illustrates one fixed-tail example:
layer pipelining, hot preloads, and seven resident hot sets reduce exposed
wait from 172 to 83, 44, and 20\,ms. It does not calibrate
(\ref{eq:depth-score}), establish multi-request WCT gains, or make seven
layers appropriate for every checkpoint state.
Here $C$ is a five-repeat median of a batch-one, 64-token native-path
forward over resident history on an A100 80GB PCIe, divided by 32 layers;
individual layer timings were not measured. The assumed bandwidth is a
25\,GB/s pool link equally shared by four recoveries, not measured H2D.
With 1\,KiB/token/layer, the Peer payload is
$7\times0.65845\times32$\,MiB $\simeq147.49$\,MiB, excluding auxiliary
state and temporary buffers. Fractions are trace-derived means, not an
integer per-layer grant. The model leaves 38.69\,ms of cold-copy work
after tail readiness; subsequent token latency requires separate evaluation.
\subsection{Adapter: Windowed Peer Capacity Control}
\label{sec:kv-runtime}
\label{sec:capacity-allocation}

Adapter (4) sets the total Peer budget $B_p$ on its target DP at each
control-window boundary. The target remains fixed within the window;
every allocation also checks physical space. The Master owns fixed-$Z$
backup decisions and per-request $L_r$ and hot sets. Adapter controls
capacity, not a common depth or an adaptive checkpoint interval.

\paragraph{Reserve memory for foreground inference.}
Subtract predicted foreground KV demand and admission headroom from the
managed KV pool, clamping the result to zero through half the pool.
The pool excludes weights and fixed runtime state but includes existing
Peer copies; $B_p$ is total permitted occupancy, not newly available space.
KV watermarks and memory-related admission stalls inform demand.
Increased pressure shrinks the budget; sustained relief allows expansion.
Report physical occupancy, ownership references, I/O leases, pending
copies, and safely reclaimable bytes separately.

\begin{figure}[t]
  \raggedright
  \refstepcounter{designalgorithm}\label{alg:adapter}
  \noindent\textbf{Algorithm \thedesignalgorithm: Windowed Peer budget update}
  \par\smallskip
  \noindent Input: pool, foreground demand, headroom, replica/lease state.
  Output: budget $B_p$, retirement work, and actual capacity reports.
  \par\smallskip
  \begin{tabularx}{\columnwidth}{@{}r@{\quad}X@{}}
    \toprule
    1 & \textbf{for each} target DP $p$ at the window boundary: \\
    2 & \quad Snapshot occupancy, owners, leases, and reservations \\
    3 & \quad Estimate foreground KV demand and headroom \\
    4 & \quad$B_p\gets\mathrm{pool}-\mathrm{demand}-\mathrm{headroom}$ \\
    5 & \quad Clamp $B_p$ to $[0,\mathrm{pool}/2]$ \\
    6 & \quad Publish $B_p$; block reservations exceeding it \\
    7 & \quad Reclaim safe zero-reference copies first \\
    8 & \quad If needed, select protected copies for demotion \\
    9 & \quad Commit DRAM; redirect; retire the old generation \\
    10 & \quad Drain old users before releasing each buffer \\
    11 & \quad Publish actual coverage and released/pending bytes \\
    \bottomrule
  \end{tabularx}
  \Description{Adapter controls only Peer capacity. When safe unreferenced
    victims are insufficient, protected copies may be demoted to committed
    DRAM, with a directory transition and I/O lease drain before release.}
\end{figure}

The migration steps are asynchronous and continue after the control tick;
a later budget update reconciles their outstanding reservations.
Algorithm~\ref{alg:adapter} does not wait for all migrations before allowing
unrelated ready requests to run.

% Capacity execution and replica lifetime, part of section 3.6.
\paragraph{Admit DRAM protection explicitly.}
Each backing pool has a finite byte ledger charging committed objects,
pending reservations, retiring buffers, and required metadata. Shared
objects are charged once. Reserve the missing dependency closure before
accepting a checkpoint; publish only completed coverage. Live checkpoint
and reader references protect DRAM objects; zero-reference objects enter
LRU reclamation, with physical release delayed until I/O leases drain.
Newer valid manifests, request completion, or an explicit transition to
recomputation release obsolete ownership. An old manifest remains retained
until its replacement is committed or its recovery claim is withdrawn.
If reservation or commit fails, keep the earlier valid frontier and replay
its missing suffix, or use R from logged inputs. A Peer denial does not
imply DRAM admission, and partial copies do not establish full D protection.
DRAM exhaustion is an operational limit, not hidden unlimited backing.

\paragraph{Protect ownership, not a permanent HBM address.}
A request reference prevents \emph{other requests} from taking its Peer
capacity, even when they have smaller slack. Normal admission uses free
space or safely reclaimable zero-reference copies; otherwise the new
Peer increment is denied. An inference request need not wait to obtain
Peer protection. The owner can replace its own hot set or reduce its
depth. Completion, including a handoff to a tool, drops that request's
references, while other owners of shared blocks retain theirs.

These ownership references differ from short-lived \emph{I/O leases}.
A lease pins the physical buffer while a kernel, read, or copy uses it.
An owner update removes only its own reference and cannot free a shared
copy still protected by another owner. Adapter contraction has explicit
system-level authority to demote even protected copies; it never overrides
an outstanding I/O lease or deletes the logical KV state.

\paragraph{Demote through a committed DRAM copy.}
The Connector and directory expose four operations:
reserve space, commit a versioned replica, acquire/release an I/O lease,
and retire a location. A block ID includes model/weights/layout, request-prefix
identity, layer, shard, and sealed token interval; a location generation
prevents reuse of a stale address. A pending copy is never readable.
A checkpoint frontier advances only when its required shards, layers,
selector state, and token log are recoverable in the declared failure model.
Independent Peer and DRAM transfers do not imply an atomic joint commit.

Retiring a Peer location follows the same order for an owner update and
an Adapter contraction:
\begin{enumerate}
  \item Reserve DRAM if needed, copy the immutable block, and commit
    a complete matching replica outside the failed source's domain.
    A valid existing DRAM replica avoids another writeback.
  \item Atomically redirect the directory to DRAM and mark the old
    Peer generation retiring. This transition closes new lease acquisition
    on that address; a lookup racing the transition must retry.
  \item Wait for all old-generation readers, writers, and DMA operations
    to release their leases. Existing holders may finish; no fresh holder
    can enter through a cached stale address.
  \item Free the Peer buffer and then return its bytes to the allocator.
\end{enumerate}
Before the directory transition, an unsuccessful writeback leaves the
surviving Peer location authoritative; afterward new readers resolve DRAM.
If DRAM allocation, transfer, or commit fails, retain the Peer copy and
report the unfulfilled shrink.
Fault detection fences stale writers and retries; directory generations
and metadata survive the source failure. These are required runtime
interfaces, not claims of a verified implementation.

\paragraph{Bound replacement and contraction.}
For old and new hot sets, retain their intersection, retire departing
owned blocks, and fetch only new arrivals. With no spare space, safe
retirement precedes admission; overlapping copies require explicit
temporary reservations. Thus changing $\{A,B,C,D\}$ to $\{C,D,E,F\}$
can retain four blocks after migration, without pinning every historical
hot set until request completion. Shared replicas and incomplete transfers
remain separately charged until physically releasable.

When occupancy exceeds a new $B_p$, stop new Peer reservations and reclaim
zero-reference copies first. If insufficient, Adapter requests protected
demotions in the Master's order: larger slack first, then deeper layer,
lower recent selection count, and stable object ID. A shared object's
priority uses the smallest slack of its remaining owners. This ordering
chooses victims; it does not serialize independent migrations or retries.
Publish the reduced coverage to the Master, and do not re-admit it beyond
the new budget. The Master subsequently reassesses request-specific $L_r$.

Contraction waits for bounded transfers and leases, not for the LLM
request to finish. Actual usage may temporarily exceed the target;
foreground admission can still wait for physical space if headroom is
exhausted. The scheduler reserves progress for required writeback, and
the controller reports pending bytes rather than claiming an instantaneous
shrink. Insufficient backing capacity remains an explicit operational limit.
\paragraph{Execute recovery with shared resources.}
Retries enter independently and advance on their own layer-readiness
events. Slack determines protection, not a request-level completion
barrier. The engine reserves a positive dispatch share for feasible recovery
work while continuing foreground batching. Within that share, ready admitted
retries receive round-robin service in bounded prefill chunks or decode steps;
eligible waiting retries use oldest-first admission subject to physical space.
Blocked layers yield rather than holding a request-level barrier. These are
required scheduler interfaces, not a claim of support in an unmodified engine.
Each retry moves through waiting-for-space, waiting-for-data, ready, and
running states. Reserve its bounded activation/KV workspace and transfer
buffers before admitting an execution chunk. The batch coordinator admits
compatible ready work and gives every TP rank the same batch, layer, and
collective order. A newly discovered miss parks that work only at a legal
kernel boundary, retaining and charging its activation buffer; it can split
from a batch only if the backend supports that split. Otherwise the whole
batch waits and the profile charges the resulting interference. Readiness
does not authorize arbitrary layer reordering or unmatched collectives.
Memory reservations precede I/O leases; release queue/directory locks before
waiting on a copy or collective, and reserve demand-transfer progress space.
Preassignment reserves neither compute nor retry workspace. Profiles
include target foreground batches, available KV/workspace, readmission
queues, shared links and failure domains. If the target has no feasible
working set, queue the retry or remap through the Router with a new
placement check; never assume that local backup alone permits execution.

Transfers contend at every shared link and at target HBM. Normal-state
backup traffic has a configured aggregate cap (default 20\% of link
capacity). During recovery or contraction, configure positive service
shares for demand fetch and required writeback while retaining a foreground
share; active shares sum to at most one per bottleneck and unused service
can be borrowed. Actual demand precedes speculative preloads within recovery.
The 20\% traffic cap is not a 20\% bound on foreground latency or overhead.
Byte-based deficit round robin across eligible requests and a bounded
outstanding-transfer count prevent one long context from monopolizing
the recovery share. Migration writeback receives its own progress share,
so shrinking does not depend on all decode requests ending. Service switches
at supported DMA/kernel boundaries, not at an assumed arbitrary instant.

The same sharing policy and foreground/retry load must enter the cost
profiles. With positive available service, finite contenders, and bounded
transfer quanta, eligible work progresses; this is not a latency guarantee
under overload. No retry must complete before another can begin merely
because its placement priority was higher.

\paragraph{Scope and validation.}
\label{sec:design-discussion}
Sparse integration needs native selection indices, recoverable selector
state and all attention branches, physical-block mapping, and the directory/
lease interface above. Poor predictions can use exact demand fetch or
full-layer loading; dense attention uses the full historical set.
Standalone requests use zero slack. Prefill failures require profiles for
the lost chunk, not a universal 64-token tail.

The present evidence consists of native selection traces, attributed
prediction results, and a single-request pipeline model using measured
forward latency and assumed transfer bandwidth. It does
not validate the depth-score profiles, failure-risk calibration, checkpoint
sensitivity, active hot-set migration, concurrent retry fairness, or
foreground interference. End-to-end tests must include recovery during
demotion, stale-location races, DRAM exhaustion, copy backlogs, and shared
bandwidth pressure. The specified interfaces and fallbacks make those
claims testable; they do not establish performance or a completed implementation.
Validation must also replay common source failures through the dependency
DAG, compare the local policy with feasible small-DAG allocations, and
report regret, including zero or negative Peer benefit. Native selector
traces and the measured batched forward do not by themselves validate
recovery equivalence to token-wise decoding.
